\clearpage

% --------------------------------------------------------------------------------------------

\section{Summary}
\label{sec:summary}

In this technical note, we have presented the results of the 3-flavour $\bar{\nu}/\nu$ joint analysis performed by the {\em VALOR} group, on the combined \RunPeriod data set. The goal of the analysis was to search for CP violation in the leptonic sector and determine limits on \deltacp as well as measure all the other oscillation parameters to which T2K is sensitive: \ssqthonethree, \ssqthtwothree, \dmsqtwothree (\dmsqonethree) and partially mass hierarchy, to which T2K has little sensitivity.

This analysis predicts
\NskOscExpectationNumuRunOneEight{} $\pm$ \NskOscExpectationNumuErrorRunOneEight{} (syst) for neutrino $\mu$-like, \NskOscExpectationNumuBarRunOneEight{} $\pm$ \NskOscExpectationNumuBarErrorRunOneEight{} (syst) for antineutrino $\mu$-like, \NskOscExpectationNueRunOneEight{} $\pm$ \NskOscExpectationNueErrorRunOneEight{} (syst) for neutrino $e$-like, \NskOscExpectationNueBarRunOneEight{} $\pm$ \NskOscExpectationNueBarErrorRunOneEight{} (syst) for antineutrino $e$-like and \NskOscExpectationNueCConepiRunOneEight{} $\pm$ \NskOscExpectationNueCConepiErrorRunOneEight{} (syst) for neutrino \nueCConepi{} events in Super K and observes respectively \NskObservedNumuRunOneEight{} for neutrino $\mu$-like, \NskObservedNumuBarRunOneEight{}  for antineutrino $\mu$-like, \NskObservedNueRunOneEight{}  for neutrino $e$-like, \NskObservedNueBarRunOneEight{}  for antineutrino $e$-like and \NskObservedNueCConepiRunOneEight{} for neutrino \nueCConepi{}. 

The $\bar{\nu}/\nu$ analysis was performed in a framework of 3-flavour oscillations including matter effects in constant-density matter. The observed reconstructed energy spectrum of single $\mu$-like ring events and reconstructed energy versus lepton angle of single $e$-like ring events in addition to $e$-like events with an additional decay electron were used to compute the confidence intervals. All \NSyst{} systematic parameters as well as the oscillation parameters which were not parameters of interest considered in this analysis were marginalised.

From the fit of the \RunPeriod data we can see that the favoured value of \deltacp is \BestFitDeltacpwRCb, close to $-\pi/2$. The Feldman-Cousins method is used to construct confidence intervals with neither over or under coverage from which the significance to exclude the hypothesis of CP conservation is greater than 2 standard deviations. 

\clearpage

%%% Local Variables:
%%% mode: latex
%%% TeX-master: "../TN327"
%%% End:
